\documentclass[12pt,a4paper]{article}
\usepackage{fontspec}
\usepackage[russian]{babel}
\setmainfont{Liberation Serif}
\usepackage{geometry}
\usepackage{booktabs}
\usepackage{amsmath}
\geometry{left=2.2cm,right=1.8cm,top=2cm,bottom=2cm}
\setlength{\parskip}{0.35em}
\setlength{\parindent}{1.25em}
\emergencystretch=3em

\begin{document}

\begin{center}
{\Large Лабораторная работа \textnumero~2}\\[0.4em]
{\large Глубокие доверительные нейронные сети}\\[0.6em]
1 октября 2026~г.
\end{center}

\section{Постановка задачи}

Нужно сжать изображения ограниченными машинами Больцмана, обученными с нуля, и оценить качество восстановления по MSE, PSNR и SSIM. Коэффициент сжатия равен 8: после кодирования представление занимает в восемь раз меньше байт, чем исходный блок.

Сравниваются три состава слоёв и три алгоритма обучения: CD-1, CD-10 и PCD.

\section{Выборка}

Взята CIFAR-10. Цветной кадр $32\times32$ переводится в яркость
\[
Y = 0{,}2989\,R + 0{,}5870\,G + 0{,}1140\,B
\]
и без перекрытия режется на блоки $8\times8$. Из одного изображения получается 16 блоков. Обучающая часть~--- все 50\,000 изображений (800\,000 блоков), тестовая~--- все 10\,000 изображений (160\,000 блоков).

Пиксели лежат в $[0,1]$. Перед обучением блоки стандартизуются по среднему и стандартному отклонению обучающих пикселей. Метрики считаются после обратного преобразования и отсечения в $[0,1]$, то есть в шкале исходной яркости.

\section{Сжатие и архитектуры}

Блок $8\times8$ в исходном виде занимает 64 байта. Скрытые нейроны бинарные, поэтому $n$ нейронов верхнего слоя~--- это $n$ бит, или $n/8$ байт. Если видимых и скрытых нейронов поровну ($n = 64$), код занимает 8 байт, и сжатие восьмикратное. Это как раз случай одной машины Гаусс--Бернулли из условия.

Промежуточное <<сжатие в $k$ раз>> для гауссовско-бернуллиевской машины понимается так же, через байты: при $k = 4$ скрытых нейронов 128, при $k = 2$ их 256. Следующие машины Бернулли--Бернулли работают уже с бинарным кодом, и двукратное сжатие у них означает уменьшение числа нейронов вдвое. У всех трёх итоговых схем верхний код состоит из 64 бит, измеренный коэффициент сжатия равен 8{,}0.

\begin{itemize}
\item Одна машина Гаусс--Бернулли, 64 видимых и 64 скрытых нейрона.
\item Гаусс--Бернулли $64 \to 128$ (4 раза по байтам) и Бернулли--Бернулли $128 \to 64$ (ещё в 2 раза).
\item Гаусс--Бернулли $64 \to 256$ (2 раза по байтам), затем Бернулли--Бернулли $256 \to 128$ и $128 \to 64$.
\end{itemize}

Дисперсия гауссовских видимых нейронов зафиксирована и равна 1, данные к этому приведены стандартизацией. Смещения и веса есть у обоих слоёв. Веса инициализируются нормальным шумом с отклонением $0{,}01$.

\section{Алгоритм обучения}

Обучение послойное и жадное. Нижняя машина видит стандартизованные пиксели. Каждая следующая получает вероятности скрытых нейронов предыдущей, а не жёсткие биты.

Контрастивная дивергенция CD-$k$ (Хинтон): положительная фаза считается по мини-выборке, скрытые вероятности не семплируются; отрицательная фаза~--- $k$ шагов выборки Гиббса, запущенных от этой же мини-выборки. В градиент отрицательной фазы входят вероятности последнего скрытого состояния. CD-1 использует один шаг, CD-10~--- десять.

PCD~--- устойчивая контрастивная дивергенция (Тилеман, 2008). Отрицательная фаза не запускается заново от данных: хранится набор частиц размера мини-выборки, и на каждом шаге он продвигается одним шагом Гиббса.

Для гауссовского видимого слоя на отрицательной фазе берётся условное среднее, а не гауссовская выборка. Скрытый слой всегда семплируется из распределения Бернулли. У бернуллиевского видимого слоя семплируются оба слоя.

Шаг обучения $10^{-3}$ для машины Гаусс--Бернулли и $10^{-2}$ для Бернулли--Бернулли, импульс $0{,}5$, затухание весов $2\cdot10^{-4}$, мини-выборка 256, 8 эпох на каждую машину.

В сам сжатый файл записывается только верхний бинарный код: единица, если вероятность скрытого нейрона не меньше $0{,}5$. Восстановление идёт сверху вниз детерминированно. Промежуточные бернуллиевские слои дают сигмоиду, нижний гауссовский слой~--- условное среднее. Затем обратная стандартизация и отсечение в $[0,1]$.

\section{Метрики}

Пусть $x$~--- исходный блок, $\hat x$~--- восстановленный, оба в $[0,1]$, $L = 1$.
\[
\mathrm{MSE} = \frac{1}{N}\sum_{i=1}^{N}(x_i-\hat x_i)^2,
\qquad
\mathrm{PSNR} = 10\log_{10}\frac{L^2}{\mathrm{MSE}_{\text{блока}}}.
\]
В таблице MSE усреднена по всем тестовым пикселям. PSNR усреднён по блокам. SSIM тоже усреднён по блокам; для блока $8\times8$ окно~--- весь блок:
\[
\mathrm{SSIM}(x,\hat x)
= \frac{(2\mu_x\mu_{\hat x}+C_1)(2\sigma_{x\hat x}+C_2)}
{(\mu_x^2+\mu_{\hat x}^2+C_1)(\sigma_x^2+\sigma_{\hat x}^2+C_2)},
\]
где $C_1 = (0{,}01)^2$ и $C_2 = (0{,}03)^2$.

Расчёт выполнен в PyTorch 2.9.1+cu128 на NVIDIA GeForce RTX 3060.

\section{Результаты}

\begin{center}
\begin{tabular}{lccc}
\toprule
Архитектура & CD-1 & CD-10 & PCD \\
\midrule
ГБ 8$\times$ & 0{,}0207 & 0{,}0185 & 0{,}0184 \\
ГБ 4$\times$ + ББ 2$\times$ & 0{,}0288 & 0{,}0266 & 0{,}0252 \\
ГБ 2$\times$ + две ББ 2$\times$ & 0{,}0364 & 0{,}0355 & 0{,}0246 \\
\bottomrule
\end{tabular}\\[0.4em]
{\small Таблица 1. MSE на тестовых блоках.}
\end{center}

\begin{center}
\begin{tabular}{lccc}
\toprule
Архитектура & CD-1 & CD-10 & PCD \\
\midrule
ГБ 8$\times$ & 17{,}74 & 18{,}21 & 18{,}22 \\
ГБ 4$\times$ + ББ 2$\times$ & 16{,}34 & 16{,}66 & 16{,}87 \\
ГБ 2$\times$ + две ББ 2$\times$ & 15{,}53 & 15{,}63 & 16{,}94 \\
\bottomrule
\end{tabular}\\[0.4em]
{\small Таблица 2. Средний PSNR тестовых блоков, дБ.}
\end{center}

\begin{center}
\begin{tabular}{lccc}
\toprule
Архитектура & CD-1 & CD-10 & PCD \\
\midrule
ГБ 8$\times$ & 0{,}5846 & 0{,}5961 & 0{,}5959 \\
ГБ 4$\times$ + ББ 2$\times$ & 0{,}4030 & 0{,}4104 & 0{,}4178 \\
ГБ 2$\times$ + две ББ 2$\times$ & 0{,}3270 & 0{,}3342 & 0{,}3557 \\
\bottomrule
\end{tabular}\\[0.4em]
{\small Таблица 3. Средний SSIM тестовых блоков.}
\end{center}

Наименьшая ошибка восстановления~--- MSE 0{,}0184 у схемы <<ГБ 8$\times$>> при алгоритме PCD. Ей соответствуют PSNR 18{,}22~дБ и SSIM 0{,}5959.

\section{Вывод}

Три требуемых состава слоёв обучены тремя алгоритмами, коэффициент сжатия во всех девяти опытах равен 8. Меньше всего ошибается одна машина Гаусс--Бернулли: MSE 0{,}0184 при PCD. CD-10 на этой же машине почти совпадает с PCD (MSE 0{,}0185), CD-1 немного хуже (0{,}0207). На двухслойной схеме порядок алгоритмов тот же, но ошибка выше: лучшая MSE равна 0{,}0252. На трёхслойной схеме PCD заметно обгоняет короткие цепи~--- 0{,}0246 против 0{,}036 у CD-1 и 0{,}035 у CD-10. При общем бюджете 64 бита дробление кода на несколько машин восстановление не улучшило.

\end{document}
